% Page 018

\seite{18}

im Seminar zu Wittlich.\ob
Die Schülerzahl in diesem Jahre war 42.\ Sämtliche\ob
Kinder sind katholisch und willfährig.

\pend
\jahresueberschrift{Das Jahr 1890}
\pstart

Mit dem Jahre 1890 begann die Schülerzahl zu steigen;\ob
sie beträgt in demselben 44 und zwar 21 Knaben\ob
und 23 Mädchen. Sämtliche Kinder sind katholisch\ob
und willfährig.

Für die Landwirte war das Jahr 1890 ein ziemlich gutes.\ob
Besonders große Erträge brachten jene Felder, welche\ob
mit sogenanntem \enquote{Kunstdünger} das bestellt\ob
waren. Dieser Edeldünger\unsicher{} ist gegenwärtig auf im\ob
ganzen Kreise durchgedrungen. Anfangs wollte die Leute\ob
diesen Dünger nicht daran; allein, es ging hiermit, wie\ob
mit so vielen Erfindungen, denn gegenüber der\ob
Thatsache Statt Widerspruch gegenüberstellen. Der\ob
Ertrag und der Preis für diesen Dünger war so groß,\ob
daß Leute und Selber, von ca 25 a groß über 300 M.\ob
\margmark{1890\\Volkszählung}%
erzielt haben.\ob
Im Jahre 1890 war 1.\ Dez. die amtliche Volkszählung:\ob
Laut dieser bestand das Dorf aus 37 Familien; dazu\ob
kommen noch 4 Familien von Klaftbrunnen, 1 von\ob
Hasselstein und 1 vom Heinzenhof. Auf Hasselstein\ob
war früher ein Stiftsgericht, in welchem sich ein\ob
Rittgutsherr V.\ befand.
\pend
